\scanpage{91}
\mspage{13}
Non seulement la condition de \deleted{\unsure{continuité}} ``prolongeabilité'' --- il vaut mieux dire d'``inductibilité'' --- sous la forme précisée, ou affaiblie (avec $\Theta$ remplacé par un $\Theta^n$), on trouve une condition supplémentaire, quand on a \inserted{de plus} un morphisme
\[
\widehat\Pi_1(C,I)\longrightarrow\widehat\Pi_1(X,I),
\]
que celui-ci passe au quotient --- ce qui s'exprime très exactement par la condition que pour \uline{un} $x_0\in I$, l'action de $\Gamma$ sur
\[
\widehat\pi_1(C,x_0)\subset\widehat\pi_1(U_{0,3},f(x_0))
\]
laisse stable le noyau de
\[
\widehat\pi_1(C,x_0)\longrightarrow\widehat\pi_1(X,x_0);
\]
alors ce sera vrai également pour tous les autres $x_i\in I$.

Il faut exprimer la possibilité de relever l'action de $\Theta$ à $\Pi_{\unsure{1|0}}(C,I)$. Notons à ce propos

\uline{Lemme} Soit $\mathcal C'-\mathcal C$ un hom de groupoïdes connexes, qui soit injectif sur les $\pi_1$, \deleted{\unsure{et}} de sorte que $\mathcal C'$ est équivalent à un revêtement de $\mathcal C$, mais on ne suppose pas que $\mathcal C'$ soit lui-$=$ un revêtement de $\mathcal C$, mais quand-même \unsure{que|pas}
\margin{90}
